\documentclass[11pt]{article}
\usepackage[margin=1in]{geometry}
\usepackage{times}
\usepackage{hyperref}
\hypersetup{colorlinks=true, linkcolor=black, urlcolor=blue, citecolor=black}
\title{The Nassikas Thruster and the Need for a Verifiable Test Record}
\author{Flyxion \\ \small Independent Researcher}
\date{}
\begin{document}
\maketitle

\begin{abstract}
The Nassikas thruster is described in the original essay collection as a superconducting resonant cavity device alleged to produce thrust through an electromagnetic interaction involving a magnetic field. The source essay offered no citations for the device's history, measured force, experimental setup, independent replication, or theoretical derivation. Searches for a traceable published paper or primary technical report were not successful in this review. The previous draft nevertheless stated specific claims about micronewton measurements, the noise floor, isolation controls, and the absence of independent replication. Those claims have been removed because they were not sourced. This revised essay preserves the original methodological thesis: a closed-system thrust claim must be evaluated by measuring external momentum transfer and excluding forces and energy flows crossing the experimental boundary. It reports no device performance that the available record does not show. The entropic-gravity framing is brief because the proposed mechanism is described as electromagnetic. This paper does not establish that the device was tested and failed; it states that the claim cannot be assessed as an experimentally demonstrated thruster without an identifiable test record.
\end{abstract}

\section{Introduction}

The original source essay attributes a reactionless-thrust proposal to Nassikas and characterizes it as involving a superconducting resonant cavity and a magnetic field. It contrasts that proposed mechanism with the better-known resonant-cavity propulsion claims associated with the EmDrive. These descriptions provide a topic for investigation, but they are not accompanied by a paper, patent, report, or other primary document that would allow the apparatus or its claims to be checked.

That evidentiary gap matters because the earlier developed version moved beyond the source record. It asserted that thrust had been measured in the micronewton range, that those readings were near the instrument noise floor, that no independent laboratory had attempted an improved-isolation replication, and that the derivation failed to conserve momentum. Without a specific source, test report, or derivation to inspect, none of those details can be responsibly represented as a documented result. This revision removes them rather than substituting a generic criticism for an actual analysis of the proposal.

\section{History and Limits of the Available Record}

The claim must be separated into three questions. First, what apparatus was proposed and what force was predicted? Second, what experiment measured the force and the system's momentum exchange? Third, what independent test, if any, reproduced the reported result? The materials identified for this review do not answer those questions. No traceable article or inventor-authored technical record was located that would verify the proposed cavity geometry, field arrangement, measurement method, numerical thrust, or test history. It is consequently not possible to determine whether the only extant descriptions are inventor-authored or whether additional primary material exists outside the sources checked.

The generic physical constraint is nevertheless clear. Internal electromagnetic forces can redistribute momentum among components of a device, but a net external impulse for the complete apparatus requires accounting for momentum transferred to something beyond the chosen system boundary. The boundary could be crossed by expelled matter, emitted radiation, interaction with an external field or structure, or an unrecognized experimental coupling. An experiment that reports a force must specify how those alternatives were measured or excluded. This does not decide the Nassikas case in advance; it defines what a persuasive test would need to establish.

Likewise, invoking superconductivity does not by itself determine the external momentum balance. Superconductors have distinctive electromagnetic properties, and their response can be analyzed within established electrodynamics. Those properties may change forces within an apparatus or produce couplings to external fields, but a particular thruster claim requires an explicit calculation and a measurement of the actual force on the entire device \cite{Jackson1999}. The name of a material or a resonant cavity is not a substitute for that calculation.

\section{What a Complete Test Must Measure}

A force sensor does not measure momentum conservation in the abstract. It measures an interaction between its sensing element and the object or support attached to it. In a small-thrust experiment, cables, thermal gradients, magnetic coupling, vibration, air currents, and electrostatic forces can all influence that interaction. These are general experimental concerns, not a claim that any one of them affected a Nassikas test. A published report would need to describe the suspension and calibration, show the sensor response with the device operating and with suitable controls, and establish that the measured force belongs to the complete apparatus rather than to a force exchanged with the support or environment.

Momentum measurements can complement a force balance. If a device is alleged to acquire an impulse, the time integral of force should correspond to a change in momentum, with the system boundary stated clearly. If energy is also claimed, the energy input and output channels must be separately recorded. A resonant cavity may store electromagnetic energy during operation, so a short observation interval should not confuse a change in stored field energy with a continuous thrust or net energy source. The relevant records would include the duration and timing of applied power, the force waveform, and the behavior after the drive is removed. None of these details can be inferred from the name of the proposed mechanism.

The distinction between a transient force and sustained propulsion is consequential. A pulse can impart momentum if the apparatus exchanges momentum with an external field, emitted radiation, or another part of the environment; observing a force during the pulse does not by itself identify that exchange. For a closed apparatus, the force and impulse must be assessed over the full operating cycle, including start-up and shut-down, rather than only during the interval that produces the most favorable sensor deflection. A reproducible signal would be an important observation, but its physical interpretation would depend on the controls and boundary accounting reported alongside it. No such signal or cycle is attributed to the Nassikas device here because no source describing one was verified.

Controls should be selected to test the mechanism rather than merely to show that a signal appears when the apparatus is energized. A powered dummy load can help identify heating or electromagnetic pickup, while a mechanically and electrically altered configuration can test whether a response follows the proposed field geometry. Reversing the device orientation can help distinguish a directional force from some balance-specific effects, although orientation tests alone do not rule out all artifacts. The methods and outcomes of these controls must be reported, not assumed. This section therefore specifies the kind of evidence that would make the claim assessable; it does not attribute any such procedure to a past Nassikas experiment.

\section{Test Record and Strongest Objection}

No published independent replication, controlled force measurement, or primary test report for the Nassikas thruster was verified in the material checked for this essay. The prior version's statements about a specific force magnitude and apparatus noise are therefore omitted. This is not a conclusion that no experiment has ever been performed. It is a conclusion about the limits of the record located here: without an identifiable report, the measurements and their controls cannot be evaluated.

The strongest objection is that the search may have missed technical material that is not indexed in common scholarly databases, and that a prototype or inventor demonstration may exist outside a journal publication. That is possible. The burden of this revision is not to infer the absence of every private or hard-to-find document. It is to distinguish the existence of a claim from evidence that the claim has been tested. If an inventor-authored source becomes available, it can establish what its author proposed or reported. It would not by itself establish an independent replication, but it would permit an accurate account of the design and test procedure.

A useful future record would identify the apparatus and its operating states, quantify the forces on the complete system, include powered and unpowered controls, monitor thermal and electromagnetic interactions with the support and sensors, and report calibration and uncertainty. Such a protocol is not presented here as a test that has already occurred. It is the minimum evidentiary structure needed to distinguish a real external impulse from forces internal to, or coupled into, the measurement system.

\section{The Entropic Gravity Framing}

The Nassikas device is presented as an electromagnetic rather than gravitational claim. The entropic-gravity framing used in this series is kept short: it treats gravitation as a large-scale emergent organization, not as a local momentum reservoir presumed available to a laboratory cavity. This framework does not decide whether a particular electromagnetic apparatus produces thrust. The question remains whether a measured external force is accompanied by a corresponding exchange of momentum across the system boundary.

\section{Conclusion}

The evidence verified for this essay does not establish the Nassikas thruster's construction, force measurements, test record, independent replication, or detailed theoretical argument. The earlier claims about specific thrust levels, noise floors, and isolation tests have therefore been removed. The general conservation-law requirement is not a substitute for examining a device-specific source; it explains why the source must document momentum transfer and experimental controls. This paper does not prove that the Nassikas thruster is impossible or that a test produced a null result. It concludes that the claim remains undocumented in the record verified here and cannot be treated as an experimentally demonstrated reactionless drive.

\begin{thebibliography}{1}

\bibitem{Jackson1999} J. D. Jackson, \emph{Classical Electrodynamics}, 3rd ed., Wiley, 1999.

\end{thebibliography}

\end{document}
